\chapter{A trainable quantum interface for ToricGT}
\label{chap:quantum-toricgt}

\reviewstatus{This chapter supplies an original, bounded implementation and
elementary proofs for a quantum feature interface to a graph transformer.
It reports nine completed GPU training runs and independent numerical checks.
It does not identify this implementation with the complete upstream ToricGT
program, a biomolecular language model, or a quantum computer. The simulation
uses three qubits per graph token. Its exact path certificates are classical
mathematical certificates, with their own proofs and verification code.}

\section{Which architecture is being extended?}

A useful connection between algebra and learning begins with a precise
interface. A theorem about a ring, a circuit, or a chain complex does not
automatically become a theorem about a trained predictor. We therefore give
each part of the construction a separate role. Graph structure specifies
legal objects; a finite circuit supplies trainable features; a transformer
predicts a value function; and an independent checker determines what a
proposed solution actually establishes.

The implementation sources reviewed for this extension are pinned to
\texttt{amelie-iska/ToricGT} commit
\texttt{a5199bc989a6fb170e21220e5e87bf95bc2fc0b4} and
\texttt{amelie-iska/ToricGT-Program} commit
\texttt{f91ba1811b5a4f70d53c7120acf8ca2a0049f85b}.
The four source modules inspected in full implement graph tokenization,
tropical attention, finite Weyl matrices, and Soft-MoE routing. The program
overview requires empirical evaluation of proposed algebraic auxiliaries;
it does not make their utility follow from their names.
\cite{qtToricGTCode,qtToricGTProgram}

Our model retains vertex and edge tokens, endpoint incidence, a mixture of
ordinary and max-plus attention, and four experts with two slots each.
Its new component is a jointly trained circuit feature encoder. The older
repository example under \path{examples/quantum_tokengt} used a frozen
three-qubit feature construction and trained a classical attention head.
The present experiment trains both the input-to-angle map and twelve circuit
offsets. This is a concrete change in the optimization problem, not an
increase in the number of qubits. The larger surface-code circuits described
elsewhere in this monograph are separate experiments.

There are also deliberate changes to the upstream graph encoding. The
reviewed tokenizer contains learned identifier embeddings. Such a table is
not automatically invariant under arbitrary renaming of vertices. Here the
identifiers enter only through equality of endpoints, equivalently through
inner products of orthonormal identifier vectors. Intrinsic byte positions
remain ordinary input attributes. A permutation of storage order transports
these attributes; it does not change the underlying string. This distinction
will be used explicitly in the equivariance proof.

The ToricBLM source located during this review is the author's June 2026
proposal, which discusses graph tokens and affine steering regions. It is
useful as a design source. An independently pinned, trained ToricBLM
repository was not located in the reviewed material. Consequently this
chapter makes no claim to reproduce its proposed 100-million-parameter
biomolecular model, its biological validation program, or broad claims of
generalization.\cite{qtToricBLM}

The selected task is exact weighted tokenization on small directed acyclic
graphs. It gives the learner a meaningful combinatorial problem while
keeping the teacher and certificate transparent. This choice also makes a
limitation explicit: dynamic programming already solves the task efficiently.
The experiment tests whether a particular trainable representation works;
it does not establish a computational advantage over that algorithm.

\section{Tokenization graphs and exact solution certificates}

\subsection{The fixed-vocabulary problem}

Let a string have length $n$. Its byte boundaries are vertices
$0,1,\ldots,n$. An allowed token occupying the substring from boundary $i$
to boundary $j$ gives an edge $e:i\to j$, where $i<j$. An edge has integer
cost $c_e$. Every one-byte edge is available, so a path to $n$ exists from
every vertex. A complete tokenization is a directed path from $0$ to $n$.
In the experiment, candidate tokens have length at most three and the
alphabet has four symbols. Costs depend on a recorded, bounded control
parameter; the full graph and its costs are supplied to the model.

The colored-DAG formulation follows the tokenization viewpoint of
ConvexTok. Its vocabulary-selection problem and its decoding problem must
be distinguished. Choosing a globally shared vocabulary under a size budget
is a discrete optimization problem. Once the allowed vocabulary is fixed,
the remaining path problem on a DAG admits the elementary solution below.
\cite{qtConvexTok}

\begin{proposition}[Bellman recursion on a finite token DAG]
\label{qt:bellman}
Suppose all edges point from smaller to larger vertex indices and every
nonterminal vertex reaches $n$. Set $d(n)=0$ and, in decreasing order of $i$,
define
\[
 d(i)=\min_{e:i\to j}\{c_e+d(j)\}.
\]
Then $d(i)$ is the minimum cost of a path from $i$ to $n$. Choosing a
minimizing edge at every visited vertex produces an optimal path.
\end{proposition}
\begin{proof}
At $n$ the only path of length zero costs zero. Suppose the assertion holds
for vertices larger than $i$. Every path beginning at $i$ has a first edge
$i\to j$ and a remaining path beginning at $j$. By the induction hypothesis
its cost is at least $c_e+d(j)$, hence at least the displayed minimum.
Conversely choose an edge attaining the minimum and append an optimal path
from its target. The result attains $d(i)$. The graph is finite and every
edge strictly increases the index, so this procedure terminates. Backward
induction proves every assertion.
\end{proof}

The recursion works even with negative edge costs because there are no
directed cycles. Our sampled costs are positive, which makes their
interpretation as token costs convenient but is not needed for this proof.
With outgoing adjacency lists, the arithmetic work is linear in the number
of edges. This is the exact teacher used during dataset generation.

\begin{proposition}[A checkable primal--dual certificate]
\label{qt:path-certificate}
Let $p:V\to\R$ satisfy $p(n)=0$ and
\[
 p(i)-p(j)\leq c_e \qquad(e:i\to j).
\]
For every path $P$ from $0$ to $n$,
$p(0)\leq C(P):=\sum_{e\in P}c_e$.
If equality holds for one legal path, that path is optimal.
The Bellman distances in Proposition~\ref{qt:bellman} give such a potential
and equality on every path obtained by choosing minimizing edges.
\end{proposition}
\begin{proof}
Write the successive vertices of $P$ as $i_0=0,i_1,\ldots,i_m=n$.
Adding its $m$ inequalities gives
\[
 \sum_{r=0}^{m-1}c_{i_r i_{r+1}}
 \geq \sum_{r=0}^{m-1}(p(i_r)-p(i_{r+1}))
 =p(0)-p(n)=p(0).
\]
Thus $p(0)$ is a lower bound for every feasible solution. A path attaining
the bound is optimal. The Bellman minimum implies
$d(i)\leq c_e+d(j)$ for each outgoing edge, and a minimizing edge makes
this inequality an equality. Summing along the resulting path proves the
last assertion.
\end{proof}

The certificate checker verifies the token substrings, continuity of the
path, its endpoint, all potential inequalities, and equality of the final
cost and lower bound. It uses literal integers, rather than a numerical
tolerance. A learned value function $\widehat d$ supplies a candidate path
by minimizing $c_e+\widehat d(j)$ at each step. Rounded predictions supply
an initial integral potential. Backward replacement
\[
 p(i)\leftarrow\min\bigl(p(i),\min_{i\to j}(c_{ij}+p(j))\bigr)
\]
makes that potential feasible while preserving $p(n)=0$. Feasibility does
not force the proposed path to attain its lower bound. This is why the
experiment reports the frequency of an immediately certified learned
proposal separately from the frequency of a merely optimal learned path.

If the certificate does not close, an exact dynamic-programming repair is
available. The final output is then certified, but its certainty belongs
to the checker and repair algorithm. Calling this a perfect neural success
rate would conceal the distinction between prediction and verification.
The stored results retain the original proposal, its regret, its initial
certificate, whether repair was required, and the final certificate.

\subsection{A separate vocabulary-selection relaxation}

For completeness, consider several token DAGs with a shared finite set of
token colors $\mathcal C$. For each document let $f_e$ indicate whether an
edge is used, and let $y_c$ indicate whether a priced token color is included.
Byte edges need no color budget. With the incidence convention
$B e=e_{\mathrm{source}(e)}-e_{\mathrm{target}(e)}$, impose
\[
 Bf=e_s-e_t,\qquad f_e\leq y_{c(e)}\text{ on priced edges},\qquad
 \sum_c y_c\leq K.
\]
The integer program uses binary variables and minimizes the number, or
specified weighted cost, of used edges. Replacing binary variables by
intervals $[0,1]$ gives a linear relaxation. A rounded vocabulary must be
evaluated by decoding again; a fractional flow is not itself a vocabulary.
Our companion implements this small formulation on the three strings
\texttt{abca}, \texttt{bcab}, and \texttt{cabc} with budget two.

\begin{lemma}[Integral unit flows on an ordered DAG]
\label{qt:unit-flow}
Let every edge of a finite graph increase a fixed vertex order. If
$f\in\N^E$ and $Bf=e_s-e_t$, with distinct vertices $s$ first and $t$ last, then $f$ is the
incidence vector of one directed $s$--$t$ path.
\end{lemma}
\begin{proof}
At $s$, incoming flow is zero and outgoing flow is one. Integrality selects
exactly one outgoing edge of multiplicity one. Follow it. At an internal
vertex reached by a unit of flow, conservation supplies an outgoing edge
with positive integral multiplicity unless the vertex is $t$. Increasing
order prevents repetition and forces eventual arrival at $t$; otherwise a
last visited nonterminal vertex would violate conservation. Subtract this
path's flow from $f$. The remainder is nonnegative and has zero boundary.
A nonzero nonnegative flow with zero boundary on a DAG is impossible:
choose its first vertex incident to a positive outgoing edge. No positive
incoming edge can enter that vertex, contradicting conservation. Therefore
the remainder is zero.
\end{proof}

\begin{proposition}[An exact lower bound from a rational LP dual]
\label{qt:lp-dual}
For the real linear program
\[
 \min c^T x,\qquad Ax=b,\quad Hx\leq h,\quad 0\leq x\leq\mathbf1,
\]
suppose $y$ is unrestricted, $z\leq0$, $u\leq0$, and
$A^T y+H^Tz+u\leq c$ coordinatewise. Then every feasible $x$ satisfies
\[
 c^Tx\geq b^Ty+h^Tz+\mathbf1^Tu.
\]
If a feasible integer solution attains this lower bound, it solves both
the relaxation and the integer problem.
\end{proposition}
\begin{proof}
Multiply the coordinatewise dual inequality by $x\geq0$ and sum. This gives
$c^Tx\geq y^TAx+z^THx+u^Tx$. The equality constraint supplies $y^Tb$.
Because $z\leq0$ and $Hx\leq h$, one has $z^THx\geq z^Th$.
Because $u\leq0$ and $x\leq\mathbf1$, one has $u^Tx\geq u^T\mathbf1$.
Combining these bounds proves the claim. Integer feasible points form a
subset of the relaxation's feasible set, so equality proves optimality in
both problems.
\end{proof}

The numerical LP solver proposes multipliers. The checker converts them
to rational numbers and explicitly checks the displayed inequalities.
If rounding a proposed multiplier changes feasibility, the coordinate
correction
$u_j=\min(0,c_j-(A^Ty+H^Tz)_j)$ restores the required inequality.
The resulting exact lower bound for the recorded example is $6$.
Exhaustive enumeration of the allowed two-color vocabularies also gives
cost $6$. The equality certifies this particular finite example. It does
not assert that the relaxation is integral for arbitrary corpora.

\section{A finite Weyl encoder with trainable gates}

\subsection{The noncommutative algebra is explicit}

Let $q\geq2$, let $\omega=\exp(2\pi i/q)$, and index the standard basis of
$\C^q$ cyclically. Define
\[
 Ue_j=\omega^j e_j,\qquad Ve_j=e_{j-1}.
\]
The backward-shift convention fixes the sign in the commutation relation.
These are the finite matrices used as an algebraic ingredient in the
reviewed ToricGT module.\cite{qtWeylCode}

\begin{proposition}[Finite Weyl relations and matrix basis]
\label{qt:weyl}
The operators $U,V$ are unitary, $U^q=V^q=I$, and $VU=\omega UV$.
The $q^2$ matrices $U^rV^s$, $0\leq r,s<q$, are an orthogonal basis of
$M_q(\C)$ for the Hilbert--Schmidt inner product.
\end{proposition}
\begin{proof}
The diagonal entries of $U$ have modulus one and $V$ permutes an
orthonormal basis, so both are unitary. Their $q$th powers are the identity.
On $e_j$ the two products satisfy
$VUe_j=\omega^j e_{j-1}$ and $UVe_j=\omega^{j-1}e_{j-1}$, proving the
relation, including the cyclic case $j=0$ because $\omega^q=1$.

If $s\not\equiv0\pmod q$, $U^rV^s$ has zero diagonal and trace zero.
If $s=0$, its trace is $\sum_{j=0}^{q-1}\omega^{rj}$, equal to zero unless
$r=0$, when it equals $q$. Indeed, for $r\neq0$ the geometric sum multiplied
by $1-\omega^r$ is $1-\omega^{rq}=0$ and $1-\omega^r\neq0$.
Using the commutation relation, the product
$(U^rV^s)^*U^{r'}V^{s'}$ is a nonzero scalar times
$U^{r'-r}V^{s'-s}$. Its trace is therefore zero for distinct pairs and is
$q$ for identical pairs. There are $q^2$ nonzero orthogonal matrices in a
$q^2$-dimensional vector space, which proves the basis assertion.
\end{proof}

This finite algebra has a perfectly ordinary classical matrix
representation. The word ``noncommutative'' establishes neither quantum
hardware usage nor a computational speedup. In the implementation $q=8$;
one vector of dimension eight is represented by three qubits, not by eight
qubits or by three independent eight-level systems. A finite representation
at a rational phase parameter also does not by itself approximate every property of an
irrational noncommutative torus.

\subsection{Circuit and measurements}

For a token feature vector $x\in\R^8$, the input encoder computes six angles
\[
 a(x)=\pi\tanh(W_a x+b_a).
\]
There are two circuit layers. In layer $\ell\in\{0,1\}$, for each
$r\in\{0,1,2\}$ apply, in this order,
\[
 R_Y(a_r+\theta_{\ell r0}),\qquad
 R_Z(a_{r+3}+\theta_{\ell r1}),
\]
to qubit $r$, then apply controlled-NOT gates $0\to1$ and $1\to2$.
The initial state is $|000\rangle$ and qubit zero is the most significant
bit. There are twelve learned offsets $\theta_{\ell rb}$. The six input
angles are reused in both layers; this matters when differentiating the
input encoder.

For the resulting unit vector $\psi(x;\theta)$ define
\[
 f(x;\theta)=
 (\Re\langle U\rangle,\Im\langle U\rangle,
  \Re\langle V\rangle,\Im\langle V\rangle,
  \Re\langle UV\rangle,\Im\langle UV\rangle),
 \quad \langle T\rangle=\psi^*T\psi.
\]
These are real features. For a unitary $T$, its real and imaginary
expectations are expectations of the Hermitian operators
$(T+T^*)/2$ and $(T-T^*)/(2i)$. Thus the feature definition has a
measurement interpretation, although estimating several noncommuting
observables generally requires separate state preparations. The present
experiments calculate expectations in a classical statevector simulator;
they do not include shot noise or noisy physical gates.

\begin{lemma}[Normalization and feature bounds]
\label{qt:bounds}
The displayed circuit preserves norm. Each of its six real features lies
in $[-1,1]$. Moreover, for unit vectors $\psi,\varphi$ and an operator $T$
with $\|T\|\leq1$,
\[
 |\psi^*T\psi-\varphi^*T\varphi|\leq2\|\psi-\varphi\|.
\]
\end{lemma}
\begin{proof}
Both Pauli rotations are unitary and controlled-NOT is a permutation of an
orthonormal basis; their product preserves norm. Cauchy--Schwarz gives
$|\psi^*T\psi|\leq\|\psi\|\,\|T\psi\|\leq1$ for the three unitary
operators under consideration. Real and imaginary parts have no larger
absolute value. Finally write the expectation difference as
\[
 (\psi-\varphi)^*T\psi+\varphi^*T(\psi-\varphi)
\]
and apply the triangle inequality and Cauchy--Schwarz to the two summands.
\end{proof}

\subsection{What the parameter-shift formula actually differentiates}

The familiar two-evaluation rule applies to an expectation whose parameter
occurs in one gate with a two-eigenvalue generator. The general analytic
gradient framework is developed by Schuld and collaborators.
\cite{QSSchuld2019} The short proof here fixes the normalization and
the composition rule needed by this implementation.

\begin{proposition}[Single-occurrence Pauli parameter shift]
\label{qt:shift}
Let $P=P^*$, $P^2=I$, and
$R(t)=\exp(-itP/2)$. For fixed state $\eta$ and Hermitian $H$, put
$F(t)=\eta^*R(t)^*H R(t)\eta$. Then
\[
 F'(t)=\tfrac12\{F(t+\pi/2)-F(t-\pi/2)\}.
\]
The same identity holds for each real feature above when the varied
offset occurs in exactly one gate of that token's circuit.
\end{proposition}
\begin{proof}
The identity $P^2=I$ splits the exponential series into even and odd powers:
$R(t)=I\cos(t/2)-iP\sin(t/2)$. On expanding $R^*HR$, the coefficients of
$H$, $PHP$, and $i(PH-HP)$ are respectively
$\cos^2(t/2)$, $\sin^2(t/2)$, and $\sin(t/2)\cos(t/2)$.
Consequently $F(t)=A+B\cos t+C\sin t$ for real constants $A,B,C$.
Its derivative is $-B\sin t+C\cos t$, and the displayed shifted difference
has exactly this value. Gates preceding $R$ can be absorbed into $\eta$;
gates following it can be absorbed by conjugating $H$. Each real or
imaginary Weyl feature uses a Hermitian observable as explained above, so
the same argument applies to each feature.
\end{proof}

\begin{proposition}[Hybrid gradient with shared parameters]
\label{qt:chain-rule}
Let $f_{tj}(\theta)$ be feature $j$ of token $t$ and let a differentiable
classical computation give loss $L=\ell((f_{tj})_{t,j})$. If offset
$\theta_k$ occurs once in each token circuit, then
\[
 \frac{\partial L}{\partial\theta_k}
 =\sum_{t,j}\frac{\partial\ell}{\partial f_{tj}}
 \frac{f_{tj}(\theta_k+\pi/2)-f_{tj}(\theta_k-\pi/2)}2,
\]
where each difference changes only that occurrence in that token circuit
and the classical derivatives are evaluated at the original features.
If a parameter occurs in several gates, sum these contributions over its
occurrences before applying its outer chain rule.
\end{proposition}
\begin{proof}
Temporarily replace every gate occurrence by an independent variable.
The ordinary multivariable chain rule expresses the derivative of $L$ as
the sum of its feature adjoints multiplied by the feature derivatives.
Proposition~\ref{qt:shift} evaluates each such derivative. Identifying
occurrence variables with the shared parameter sums their derivatives,
again by the chain rule. This proves both statements.
\end{proof}

It would generally be incorrect to replace the derivative of the entire
nonlinear loss by half the difference of two shifted losses. For example,
if $f(t)=\cos t$ and $L(t)=f(t)^2$, those two shifted losses are equal,
whereas $L'(t)=-\sin(2t)$ is usually nonzero. The graph attention and
expert networks are nonlinear functions of the features. Our verification
therefore compares their feature adjoints composed with independent
parameter-shift Jacobians against automatic differentiation of the full
loss. At max-plus ties a classical derivative may fail to exist; the
implementation uses its framework's subgradient convention. The analytic
identity above is asserted at differentiability points, not as a unique
derivative at every tie.

The circuit offsets each have a single occurrence per token. An input
angle occurs once in each of the two layers, so its derivative has two
occurrence contributions. The input-to-angle map then contributes its
ordinary derivative $\pi\operatorname{sech}^2(W_ax+b_a)$. Automatic
differentiation performs this composition during training; the independent
shift check specifically audits the shared circuit offsets.

\section{A graph transformer whose symmetries can be stated precisely}

\subsection{Tokens and incidence}

For a vertex $i$, use endpoint pair $(i,i)$; for an edge $i\to j$, use
$(i,j)$. Each token has eight scalar attributes: vertex and edge indicators,
the two byte positions divided by twelve, token span divided by twelve,
cost divided by thirty-two, the bounded control divided by two, and the
edge's cost slope divided by two. Inapplicable entries are zero.
The endpoint-incidence matrix is
\[
 R_{tu}=\tfrac14\sum_{a,b=1}^2
   \mathbf1\{\operatorname{end}_a(t)=\operatorname{end}_b(u)\}.
\]
Equivalently replace the indicator by the inner product of orthonormal
endpoint vectors. This encoding uses equality structure, not a coordinate
choice in the identifier space. Vertex and edge tokenization has a primary
antecedent in TokenGT; its general expressivity results involve additional
architecture hypotheses and are not claimed for this small model.
\cite{qtTokenGT}

The initial hidden table is
\[
 X_t=W_x x_t+b_x+W_f f(x_t;\theta)+b_f\in\R^{24}.
\]
It has two graph blocks and two attention heads per block. Each block
normalizes token rows, forms learned queries, keys, and values, and centers
and normalizes each query and key before forming scores. For one head the
scores have the form
\[
 S_{tu}=\frac{q_t^T k_u}{\sqrt{12}}+\rho R_{tu},\qquad \rho>0.
\]
The two aggregation operations are
\[
 A^{\rm soft}_{t r}=\sum_u
 \frac{e^{S_{tu}}}{\sum_v e^{S_{tv}}}v_{ur},
 \qquad
 A^{\rm max}_{t r}=\max_u(S_{tu}+v_{ur}).
\]
A learned sigmoid gate mixes them. The head outputs are concatenated and
mapped back to width twenty-four with a residual connection. The graph
here is fully observed; no autoregressive next-byte probability is being
reported. A causal language model would require a separate reveal-order
mask and proof that endpoint attributes expose no future information.

\begin{lemma}[Associativity of finite max-plus contraction]
\label{qt:maxplus}
For finite compatible matrices with entries in
$\R\cup\{-\infty\}$, define
$(A\odot B)_{ik}=\max_j(A_{ij}+B_{jk})$, with the usual convention for
$-\infty$ and $\max\varnothing=-\infty$. Then $(A\odot B)\odot C=A\odot(B\odot C)$.
Splitting the contracted index set into blocks and taking the maximum of
the block maxima gives the same result.
\end{lemma}
\begin{proof}
Both parenthesizations at entry $(i,l)$ equal
$\max_{j,k}(A_{ij}+B_{jk}+C_{kl})$. Finite maxima can be taken in either
order. A maximum over a finite disjoint union is the maximum of the maxima
over its pieces, proving the second statement.
\end{proof}

This justifies an exact blockwise max-plus reduction when its scores and
values are already fixed. It does not make the whole graph network
piecewise affine in its raw input or in its parameters: normalized
queries and keys, their dot products, quantum expectations, softmax,
layer normalization, and GELU activations are nonlinear.

\subsection{Soft experts and a permutation proof}

For an $m\times d$ token table $X$ and eight learned slot keys, form logits
$L_{ts}$ by scaled cosine similarity. In the implementation each vector
$v$ is normalized as $v/\max(\|v\|_2,10^{-6})$, including the zero vector;
this convention is also used for the centered queries and keys. Row layer
normalization uses its population variance plus $10^{-5}$ under the square
root. These positive constants make every displayed row operation defined
at constant or zero rows. Define two differently normalized
arrays
\[
 D_{ts}=\frac{e^{L_{ts}}}{\sum_{u=1}^m e^{L_{us}}},\qquad
 C_{ts}=\frac{e^{L_{ts}}}{\sum_{r=1}^8 e^{L_{tr}}}.
\]
The slot inputs are $\widetilde X=D^TX$. Two slots are assigned to each
of four learned feed-forward experts, giving output slots
$\widetilde Y$. The token output is $C\widetilde Y$, multiplied by a
learned residual coefficient. Dispatch normalizes over tokens; combination
normalizes over slots. Confusing those axes changes the model.
\cite{qtSoftMoE}

\begin{proposition}[Equivariance of the implemented mathematical architecture]
\label{qt:equivariance}
Let $P$ be a permutation matrix on the valid tokens of a nonempty graph.
Transport token attributes and incidence by $X\mapsto PX$ and
$R\mapsto PRP^T$. Every displayed graph block and the rowwise readout
commute with this action. Thus the predicted node values are permuted
along with the corresponding node tokens. A coordinated renaming of
endpoint identifiers also leaves the incidence data unchanged.
\end{proposition}
\begin{proof}
Rowwise affine maps, rowwise normalizations, coordinatewise activations,
and the identical token circuit commute with left multiplication by $P$.
Queries, keys, and values therefore transform to $PQ,PK,PV$, so their
score matrix transforms to $PSP^T$. Row softmax transforms by the same
two-sided permutation, and its product with the permuted values is the
permutation of the original output. For max-plus attention, relabeling
the finite maximizing index merely reorders the same candidates, giving
the same conclusion. Residual addition and concatenation preserve it.

For the expert layer, $L$ transforms to $PL$. Consequently dispatch and
combination transform to $PD$ and $PC$. The slot inputs are unchanged:
$(PD)^T(PX)=D^TX$. Each expert acts on those unchanged slots, so the
output slots are unchanged. Multiplication by $PC$ permutes the final
token outputs. Induction over the blocks and the final rowwise readout
proves the assertion. Finally, a bijective renaming preserves every
endpoint equality appearing in $R$, which proves the last statement.
\end{proof}

The claim concerns attributed graphs. Changing a byte position while
keeping a token in place changes an attribute and can change the answer.
Likewise, erasing edge direction or swapping the semantic source and
target is not the token-storage permutation in the proposition. Masks
must be transported too when padding or causal restrictions are present.
In exact mathematics, padded keys are absent from sums and maxima. The
implementation checks the relevant finite-precision masking behavior
numerically, including insertion of padding and token permutations.

\section{Toric path relations, graph rewrites, and steering regions}

\subsection{A small commutative algebra with an exact graph meaning}

Let $G$ be a finite directed graph and give its edge variables the incidence
degrees $d_e=e_{\mathrm{target}(e)}-e_{\mathrm{source}(e)}\in\Z^V$,
the columns of a matrix $D$. Thus $D=-B$ for the flow matrix above, which does not
change the kernel. Consider the algebra map
\[
 \Phi:k[z_e:e\in E]\longrightarrow
 k[t_v^{\pm1}:v\in V],\qquad
 z_e\longmapsto t_{\mathrm{target}(e)}/t_{\mathrm{source}(e)}.
\]
For a nonnegative exponent vector $u$, its image has exponent vector
$Du$ with this incidence convention. Multiplication forgets the order
in which edges were listed; it retains their net boundary.

\begin{proposition}[The binomial kernel of the boundary code]
\label{qt:binomial-kernel}
The kernel of $\Phi$ is the ideal generated by all
$z^u-z^v$ with $u,v\in\N^E$ and $Du=Dv$.
In particular, two directed paths with the same initial and final vertices
have the same image under $\Phi$.
\end{proposition}
\begin{proof}
Every displayed binomial maps to zero. Conversely write a polynomial in
the kernel as a finite sum $\sum_u c_u z^u$ and group its terms according
to the integer vector $Du$. Distinct Laurent monomials are linearly
independent over $k$. Thus within each group the coefficients sum to zero.
Choose one exponent $u_0$ in a nonempty group. That group's contribution
equals $\sum_{u\ne u_0}c_u(z^u-z^{u_0})$, a sum of the stated binomials.
Adding the finitely many groups proves the kernel assertion. Along a path,
the incidence columns telescope to $e_t-e_s$, proving the final claim.
\end{proof}

For three arrows $a:0\to1$, $b:1\to2$, and $c:0\to2$, the two complete
paths give the relation $z_a z_b-z_c$. In this example the whole kernel is
$(z_a z_b-z_c)$. One may eliminate $z_c$, leaving the polynomial ring in
$z_a,z_b$; their incidence vectors are linearly independent, so their
images satisfy no further polynomial relation. The unit-flow polytope is
\[
 (f_a,f_b,f_c)=(\lambda,\lambda,1-\lambda),\qquad0\leq\lambda\leq1.
\]
The two integer endpoints are the two tokenizations. For edge costs
$c_a,c_b,c_c$, the preferred endpoint changes on the wall
$c_a+c_b=c_c$. The same comparison determines which side of the path
binomial is selected by the corresponding weight. This is an explicit
connection among a graph problem, a polytope, a binomial ideal, and a
piecewise-linear optimal value function.

The kernel relation alone does not certify an arbitrary sequence of
rewrite operations. The vector $u-v$ may have negative entries, and
applying it to a current state may leave the nonnegative orthant.
Even nonnegative boundary-preserving changes on a cyclic graph can add
a disconnected circulation. The DAG and unit-flow hypotheses in
Lemma~\ref{qt:unit-flow} are what turn a nonnegative solution with the
correct boundary into one legal path.

\begin{proposition}[A constrained rewrite preserves legal tokenization]
\label{qt:rewrite}
On the ordered DAG of Lemma~\ref{qt:unit-flow}, let $x$ be a legal path
flow and let $\delta\in\Z^E$ satisfy $B\delta=0$.
If $x+\delta\geq0$ coordinatewise, then $x+\delta$ is another legal
path flow with the same endpoints. Its cost change is exactly $c^T\delta$.
\end{proposition}
\begin{proof}
Linearity gives $B(x+\delta)=Bx+B\delta=e_s-e_t$.
The assumed nonnegative integral vector therefore satisfies
Lemma~\ref{qt:unit-flow} and is a path. The cost identity follows by
distributivity: $c^T(x+\delta)-c^Tx=c^T\delta$.
\end{proof}

A graph grammar can use such a local replacement when its interface and
nonnegativity tests are explicit. If declared path relations are made
columns of a map $\partial_2:C_2\to C_1$, their zero boundary gives
$\partial_1\partial_2=0$. A replacement by $\partial_2 w$ preserves the
homology class of a cycle. This is a genuine chain-complex connection to
the earlier chapters. It is weaker than proving that every algebraically
equivalent trajectory is semantically interchangeable. Dependencies,
labels, resource limits, and external facts need their own preservation
conditions. The commutative algebra records the chosen invariant; it
cannot reconstruct information intentionally discarded by abelianization.

\subsection{Affine margins and nonlinear circuit coordinates}

Suppose a finite decision system uses affine scores
$\ell_j(z)=a_j^Tz+b_j$ on a feature space. A fixed strict sign pattern of
these scores determines an open polyhedral region. It may be empty, and
there may be many disconnected regions after pulling it back through a
nonlinear encoder. The affine-space statement is the useful part of the
steering proposal; identifying a region with a scientifically correct
reasoning mechanism would require additional evidence.

\begin{proposition}[Preserving a finite affine sign pattern]
\label{qt:margin}
Let $s_j\in\{-1,1\}$ and suppose
$s_j\ell_j(z)\geq\gamma>0$ for every $j$. If
$|a_j^T\Delta z|<\gamma$ for every $j$, then
$s_j\ell_j(z+\Delta z)>0$ for every $j$.
In particular it suffices that
$\|\Delta z\|_2<\gamma/\max_j\|a_j\|_2$ when the denominator is positive.
If every $a_j=0$, every perturbation preserves the signs.
\end{proposition}
\begin{proof}
For each $j$,
$s_j\ell_j(z+\Delta z)=s_j\ell_j(z)+s_j a_j^T\Delta z
\geq\gamma-|a_j^T\Delta z|>0$.
Cauchy--Schwarz gives the sufficient norm bound. If all coefficients
vanish, all scores are constant.
\end{proof}

The proof applies, for example, to a finite list of path-cost comparisons
when the edge costs depend affinely on $z$. To use it in circuit parameters
one also needs a bound on the feature change. Here is a directly proved
bound for the twelve independent offsets, with the token input held fixed.

\begin{lemma}[A conservative circuit-offset perturbation bound]
\label{qt:circuit-lipschitz}
For the token circuit above and two offset vectors $\theta,\theta'$,
\[
 \|\psi(\theta')-\psi(\theta)\|
 \leq\tfrac12\|\theta'-\theta\|_1,
 \qquad
 \|f(\theta')-f(\theta)\|_2
 \leq\sqrt6\,\|\theta'-\theta\|_1.
\]
\end{lemma}
\begin{proof}
For a Pauli rotation the derivative has operator norm $1/2$.
Integrating along the real angle interval gives
$\|R(t')-R(t)\|\leq|t'-t|/2$.
Subtract two products of unitary gates by replacing one factor at a time.
The operator norm of each difference term is at most the norm difference
of that one factor, because all other factors are unitary. The sum over
the twelve independently varied gates proves the state bound.
Lemma~\ref{qt:bounds} then bounds the absolute change of each of the six
features by $\|\theta'-\theta\|_1$. Summing their squared bounds proves
the second inequality.
\end{proof}

Combining these propositions gives a sufficient margin condition for an
\emph{affine readout directly on the six features}. It does not give the
same constant for the full transformer. The latter's nonlinear layers
and learned weights contribute further Lipschitz factors; max-plus ties
also require the appropriate nonsmooth interpretation. Moreover a
polyhedron in expectation-feature space need not have a polyhedral
preimage in angle space. Even the one-qubit feature $\cos\theta$ has
periodic, nonlinear level sets. This qualification prevents a useful
local sign guarantee from becoming an unsupported global geometry claim.

\section{Three source corrections with complete small counterarguments}

The reviewed ToricGT V text contains several elementary interfaces that
are useful when their hypotheses are retained. The following localized
qualifications concern its saturation example, a sign in its predictive
transfer definition, and a proposed monotonicity interpretation of
coverage homology. The original pinned sources are preserved. The review
ledger records the exact commit, path, and line ranges; the corrections
here are new mathematical explanations, not silent edits to that source.
\cite{qtToricGTAlgebra}

\subsection{Saturation is relative to the generated lattice}

\begin{proposition}[The parity semigroup is normal in its own group]
\label{qt:parity-normal}
Let
$Q=\N(2,0)+\N(0,2)+\N(1,1)\subseteq\Z^2$. Then
\[
 \operatorname{gp}(Q)=\{(a,b)\in\Z^2:a+b\text{ is even}\},
 \qquad
 Q=\operatorname{gp}(Q)\cap\R_{\geq0}^2.
\]
In particular $Q$ is saturated in $\operatorname{gp}(Q)$, and $(1,0)$ is
not a saturation hole in that group.
\end{proposition}
\begin{proof}
Each generator has even coordinate sum, so every integer combination
does too. Conversely if $a+b$ is even, then $a-b$ is even and
$(a,b)=((a-b)/2)(2,0)+b(1,1)$ is an integer combination of generators.
This proves the group description.

For nonnegative $(a,b)$ with even sum, the coordinates have the same
parity. If both are even, use $(a/2)(2,0)+(b/2)(0,2)$. If both are odd,
use $((a-1)/2)(2,0)+((b-1)/2)(0,2)+(1,1)$. All coefficients are
nonnegative integers, proving the intersection formula. If
$m(a,b)\in Q$ with $m>0$ and $(a,b)\in\operatorname{gp}(Q)$, then
$a,b\geq0$, so the formula implies $(a,b)\in Q$.
Finally $(1,0)$ has odd coordinate sum and is outside the generated group.
\end{proof}

The fact that $2(1,0)\in Q$ does not contradict normality: its antecedent
omits membership in $\operatorname{gp}(Q)$. If one instead studies the
larger ambient lattice $\Z^2$, that is a different lattice specification.
This is the same distinction that matters when proving normality of the
semigroup degenerations in the standard-monomial chapters.

\subsection{Inclusions do not make Betti numbers monotone}

\begin{proposition}[A point, a circle, and a filled triangle]
\label{qt:betti-counterexample}
There exist inclusions of finite simplicial complexes
$K_0\subset K_1\subset K_2$ for which the first Betti numbers over
$\Q$ are $0,1,0$, respectively. Thus inclusion alone implies neither
nondecreasing nor nonincreasing first Betti numbers.
\end{proposition}
\begin{proof}
Take $K_0$ to be vertex $0$, $K_1$ the boundary of the triangle on
vertices $0,1,2$, and $K_2$ the filled triangle. In $K_1$, orient the
edges $0\to1$, $1\to2$, $2\to0$. The boundary matrix is
\[
 \partial_1=
 \begin{pmatrix}-1&0&1\\1&-1&0\\0&1&-1\end{pmatrix}.
\]
Its first two columns are independent and its three columns sum to zero,
so its rank is two and its kernel is spanned by $(1,1,1)^T$.
There are no two-simplices in $K_1$, giving $\dim H_1(K_1;\Q)=1$.
In $K_2$, orient the added triangle so that
$\partial_2(1)=(1,1,1)^T$. This map has rank one and its image equals
$\ker\partial_1$; hence $H_1(K_2;\Q)=0$.
The point has no one-chains, so $H_1(K_0;\Q)=0$ as well.
\end{proof}

A growing coverage complex can create a class and later kill it. A
statement that a support construction grows under evidence addition may
be correct while a monotone ``hole penalty'' conclusion is not justified.
One must specify induced homology maps, persistence intervals, and what
the complexes represent. Even a genuine persistent class is a property
of that representation, not a theorem that a scientific assertion is
false or that all relevant evidence has been gathered.

\subsection{Use one sign for predictive gain}

For normalized predictive distributions with positive probabilities on
the observed symbols, define the unaugmented codelengths
\[
 C_0=-\sum_t\log_2q_0(b_t\mid b_{<t}),\qquad
 C_1=-\sum_t\log_2q_1(b_t\mid b_{<t}).
\]
Their gain is
\[
 C_0-C_1=\sum_t\log_2\frac{q_1(b_t\mid b_{<t})}
                                  {q_0(b_t\mid b_{<t})}.
\]
If a deployed auxiliary costs $B$ additional bits, the total codelength
gain is $C_0-C_1-B$. Positive gain means improvement with this convention.
For a single observed symbol with $q_0=1/4$ and $q_1=1/2$, the
unaugmented gain is one bit. Reversing the ratio gives minus one bit.
Artifact-adjusted gain and ordinary predictive bits per byte should be
named separately, rather than changing the metric's definition mid-proof.
These identities say nothing about the task accuracy of a neural model
unless the probabilities and evaluation protocol have been supplied.

The companion's \path{algebra_checks.py} verifies the displayed boundary
ranks and zero composition exactly over $\Q$, checks the parity
decomposition on $545$ bounded pairs, and records the rational probability
example. The general semigroup claim follows from its proof above, not
from finite enumeration. The finite checks provide reproducible
illustrations and catch transcription mistakes.

\section{Completed learning experiment and numerical evidence}

\subsection{Protocol and comparators}

The fixed dataset contains $512$ training graphs and $128$ graphs in each
of the validation, test, and length-shift sets. All underlying strings
are distinct across the four sets. Training, validation, and ordinary
test strings have lengths four through eight; the length-shift set has
lengths nine through twelve. Token availability and integer cost parameters
are sampled from the recorded generator. Each byte fallback costs sixteen
integer units. In the following tables one unit of reported cost equals
that byte-fallback cost, so integer teacher costs are divided by sixteen.
The network itself predicts distances divided by thirty-two.

All three model variants use the same graph encoder layout, width
twenty-four, two attention blocks, four experts per block, and scalar node
readout. The trainable quantum and classical variants each contain
$25{,}515$ trainable parameters. The classical feature comparator uses
the same six encoded angles and twelve phase parameters, with bounded
features
\[
 g_j(a;p)=\tfrac12\bigl(\sin(a_j+p_j)
             +\cos(a_{j-1}+p_{j+6})\bigr),\qquad j\pmod6.
\]
The other input and graph layers are identical. This is a matched
parameter-count comparator, not a claim that it exhausts possible
classical architectures. In the frozen-circuit ablation, both the
$54$-parameter angle encoder and twelve circuit offsets are frozen, while
the remaining $25{,}449$ parameters are trained.

Each variant is trained at seeds $17,29,43$ for eighty epochs, batch size
thirty-two, AdamW learning rate $0.0015$, weight decay $10^{-4}$, and a
cosine learning-rate schedule ending at one tenth of its initial value.
The loss is mean squared error over valid node predictions, with padded
and edge-token targets excluded. The same initialization seed and
minibatch ordering are used for the corresponding variants. For each run,
the checkpoint with lowest validation node loss is selected before its
test and length-shift metrics are computed. The representative browser
checkpoint is the quantum run with lowest validation loss, not a
checkpoint chosen for its test score. A preceding two-epoch smoke run
checked execution and export; its results were not candidate checkpoints
for the reported experiment.

The runs used PyTorch $2.8.0$ with CUDA $12.8$ on an NVIDIA GeForce RTX
4090. They were classical GPU statevector simulations with analytic
expectations and automatic differentiation, not quantum-hardware
experiments and not cuStateVec training. The existing environment was
used without changing its packages. Process allocation was capped at
one tenth of GPU memory. The largest measured PyTorch allocation during
the nine runs was $183{,}426{,}048$ bytes, about $175$ MiB; this allocator
measurement does not include every CUDA-context allocation or the memory
of other applications. The full set of runs completed in approximately
$109$ seconds.\cite{qtExperiment}

\subsection{Results, including the negative evidence}

Table~\ref{qt:benchmark} reports means and sample standard deviations over
the three seeds. ``Optimal'' refers to the original learned proposal.
The initial certificate is stricter: the proposal must attain the lower
bound produced from its rounded, repaired potential. A proposal may be
optimal even when that particular potential does not certify it.
All final paths pass exact certification after the separately reported
dynamic-programming repair where required.

\begin{table}[htbp]
\centering\small
\begin{tabular}{llrrr}
\toprule
Set & Variant & Node RMSE & Optimal (\%) & Initial cert. (\%)\\
\midrule
Test & Trainable quantum & $0.265\pm0.003$ & $83.85\pm2.39$ & $24.74\pm0.45$\\
Test & Classical & $0.266\pm0.009$ & $84.38\pm0.78$ & $30.47\pm5.63$\\
Test & Frozen quantum & $0.284\pm0.011$ & $83.07\pm1.97$ & $25.78\pm5.12$\\
\midrule
Longer & Trainable quantum & $0.684\pm0.150$ & $54.69\pm3.41$ & $6.51\pm3.85$\\
Longer & Classical & $0.790\pm0.117$ & $50.00\pm4.75$ & $5.47\pm3.41$\\
Longer & Frozen quantum & $0.676\pm0.069$ & $52.86\pm0.90$ & $9.90\pm2.26$\\
\bottomrule
\end{tabular}
\caption{Measured learning outcomes on the fixed token-DAG dataset.
Uncertainty entries are sample standard deviations over three seeds,
not confidence intervals or statistical significance claims.}
\label{qt:benchmark}
\end{table}

The mean proposal cost regret on the ordinary test set is $0.0386$ for
the trainable quantum model, $0.0356$ for the classical comparator, and
$0.0399$ for frozen quantum features. On longer strings these means are
$0.2246$, $0.2625$, and $0.2324$, respectively. The mixed ordering is
substantive evidence: the classical comparator has slightly better
ordinary-test path accuracy and regret, while the trainable quantum model
has somewhat better mean path accuracy and regret under this length shift.
Frozen quantum features have a slightly lower mean length-shift node
RMSE than the trainable circuit. Three seeds on one synthetic task do not
support a general superiority conclusion for any of the variants.

The twelve quantum offsets change by Euclidean norms $0.146$, $0.151$,
and $0.158$ at the selected checkpoints, while the corresponding changes
in the frozen ablation are exactly zero. Initial gradients for the
circuit offsets and angle encoder are recorded. This establishes that
the new circuit feature map participates in training. It does not show
that every parameter is identifiable, that all gradients are equally
useful, or that optimization found a global minimum.

A batch of thirty-two graph forwards takes approximately $2.56$ ms for
the trainable quantum variant and $1.81$ ms for the classical comparator
on the measured GPU, after warm-up and synchronization. These are small
model-forward timings, excluding exact path repair and data preparation.
The exact dynamic program is an algorithmic baseline with zero regret;
the model comparison does not establish an end-to-end speed advantage
over it. No inference about a physical quantum device follows from these
classical timings.

\subsection{Independent checks and browser execution}

The circuit is implemented three ways: indexed PyTorch tensor operations,
independent dense NumPy Kronecker matrices, and a dependency-free JavaScript
statevector reference. The GPU verification uses double precision and
reports a maximum PyTorch--NumPy state difference of $3.73\cdot10^{-16}$,
feature difference of $4.45\cdot10^{-16}$, and automatic-differentiation
versus parameter-shift Jacobian difference of $3.34\cdot10^{-16}$ on
the recorded cases. The full nonlinear-loss chain-rule check differs by
$1.12\cdot10^{-16}$. These are observed numerical agreement bounds on
finite tests, not formal floating-point error theorems.

The same verification includes forty-eight independent exhaustive path
enumerations, exact rational LP-dual checks, token permutation and
padding checks, and the two normalization axes of the expert router.
Finite differences provide an additional numerical gradient check, but
the analytic parameter-shift identity is the mathematical reference.
All source and dataset hashes associated with the training runs are
retained. Checkpoints use ordinary NumPy array archives, and per-example
predictions and certificates are stored in JSON. The complete results
of all nine runs are retained, including the runs that do not favor
the quantum variant.

The browser runs the whole selected graph model. Its
\texttt{runGraph} function accepts the graph, exported weights and optional
circuit offsets. It recomputes input
angles, circuit expectations, both attention blocks, expert routing,
node predictions, and the candidate path. The optional offset tensor
has shape $2\times3\times2$. It changes the circuit offsets for every
token and propagates the resulting features through the full model.
The interface then checks the candidate and, if necessary, displays the
separate exact repair. Recorded predictions are not substituted for the
live computation.

JavaScript uses double arithmetic and an approximation to the error
function in GELU. Against saved GPU predictions on two ordinary-test and
two length-shift graphs, the largest discrepancy in normalized node
prediction is $6.81\cdot10^{-7}$. The four learned paths and final exact
certificates agree. Increasing one circuit offset by $0.4$ radians changes
the full model's predictions by as much as $0.555$ in the checked example;
the interactive control therefore changes the actual model computation.
This comparison is recorded with the hashes of the JavaScript modules
and exported weights.\cite{qtExperiment}

\section{Connections that follow, and research questions that remain}

The finite Weyl algebra is the strongest direct connection to quantum
algorithm theory in this experiment. Clock and shift matrices encode
finite translation and character operations. Similar algebra is central
to finite Fourier analysis. Here, however, their expectations are learned
features of a very small circuit. No hidden subgroup is identified, no
elliptic-curve discrete logarithm is extracted, and no improvement to
Shor's complexity follows. The earlier cryptographic chapters identify
the additional problem-specific structures required for those claims.

The binomial path kernel connects more directly to the standard-monomial
and semigroup chapters. It is an actual quotient map with a proved kernel,
and its simplest path comparison is an actual initial-term wall. The
parity example shows why the lattice used for normalization must be
specified. A hidden feature called a ``toric coordinate'' is not by
itself a point on a chosen toric variety; one must supply the monomial
map, relations, and any positivity or saturation hypotheses. Our finite
path construction supplies those data in an elementary setting.

Graph rewrites connect to complexes through the condition
$\partial_1\partial_2=0$ and to learning through proposal and acceptance.
A trainable model may suggest a boundary-preserving replacement; exact
checks enforce the interface and nonnegativity constraints. This is a
practical division of labor for a graph grammar. Extending it to arbitrary
typed hypergraphs requires the appropriate matching and gluing conditions,
which are developed separately in the graph-rewrite chapter. It is not
justified by a count-vector identity alone.

There is also a useful possible quantum extension whose theorem can be
stated independently of its implementation. It makes precise what it
would mean for a variational state to remain within a declared feasible
space. It is not the circuit trained in the preceding experiment.

\begin{proposition}[A constraint-preserving variational interface]
\label{qt:feasible-hilbert}
Let $\Omega$ be a finite set, $F\subseteq\Omega$ nonempty, and
$\Pi_F$ the orthogonal projector onto the span of basis vectors
$|x\rangle$ with $x\in F$. Let $\psi_0$ be a unit vector in this subspace.
If every gate $G_j$ is unitary and commutes with $\Pi_F$, the resulting
state $\psi=G_m\cdots G_1\psi_0$ remains in that subspace. Measurement in
the indicated basis returns an element of $F$ with probability one.
For a real cost function $c$ on $F$, put
$H_c=\sum_{x\in F}c(x)|x\rangle\langle x|$ on this subspace and
$c_* =\min_{x\in F}c(x)$. Then
$\langle H_c\rangle\geq c_*$. If every nonoptimal cost is at least
$c_*+\Delta$ for some $\Delta>0$, the probability of a nonoptimal outcome
is at most $(\langle H_c\rangle-c_*)/\Delta$.
\end{proposition}
\begin{proof}
Commutation gives
$\Pi_F G_j\eta=G_j\Pi_F\eta=G_j\eta$ whenever $\Pi_F\eta=\eta$.
Induction proves subspace preservation, and unitarity preserves norm.
Writing $\psi=\sum_{x\in F}\alpha_x|x\rangle$ gives
$\sum_x|\alpha_x|^2=1$ and
$\langle H_c\rangle=\sum_x c(x)|\alpha_x|^2$.
This is a convex combination of costs and is at least $c_*$.
Subtracting $c_*$ and retaining only nonoptimal terms yields
\[
 \langle H_c\rangle-c_*
 \geq\Delta\sum_{x:c(x)>c_*}|\alpha_x|^2,
\]
which is the claimed probability bound.
\end{proof}

To make this a useful algorithm one must efficiently encode $F$, prepare
$\psi_0$, construct expressive gates commuting with its projector, and
estimate the cost. Those are substantial additional tasks. In particular,
the three-qubit feature circuit in this chapter has no basis state for
every tokenization path of an arbitrary graph. Its unitarity certifies
normalization, while the separate graph checker certifies tokenization.
Confusing these two invariants would erase the main proof boundary.

Finally, the finite graph and tensor operations do not by themselves
establish a continuum gauge theory, an elliptic index theorem, a positive
Yang--Mills mass gap, or a cosmological model. They can provide finite
data structures and exactly checked algebraic interfaces for such
research. A field-theoretic or cosmological conclusion additionally
requires the physical interpretation, limiting procedure, analytic
hypotheses, and observational tests appropriate to that problem. The
cosmology and gauge-theory chapters retain those requirements. The
achievement established here is narrower and concrete: a source-grounded
trainable quantum graph feature model, explicit mathematical guarantees
for its finite interfaces, measured comparator results, and a live
reproducible implementation.


